\node[tag,right=2mm of isa.east,anchor=west] {\S\ref{sec:isa}};
\node[tag,right=2mm of frame.east,anchor=west] {\S\ref{sec:hardware}};
\node[tag,right=2mm of geo.east,anchor=west] {\S\ref{sec:geometry}};
\node[tag,right=2mm of phys.east,anchor=west] {\S\ref{sec:physics}};

\begin{scope}[on background layer]
  \node[draw=rule,dashed,rounded corners=4pt,fit=(net)(phys),inner sep=7pt] {};
\end{scope}
\end{tikzpicture}
\end{center}

\noindent The unusual property of this stack is that the layers are not independent
design choices stacked by convention. Each one is \emph{forced} by the one below it. The
network is $40$-ary because the geometry has $40$ points. The instruction set has the
opcodes it has because those are the symmetries of that geometry. This is what the rest
of the document demonstrates, layer by layer.

% =====================================================================================
\section{Qutrits, and what ``self-entangled'' means}\label{sec:physics}
% =====================================================================================

\subsection{From bits to qutrits}

\begin{plain}
A bit is a switch: off or on, $0$ or $1$. A \emph{trit} is a three-way switch: $0$, $1$,
or $2$. Ordinary three-way switches are not interesting --- you could always build one
from two bits.

A \emph{qutrit} is different. It is a physical three-way switch obeying quantum
mechanics, which means it can be in a blend of all three settings at once, and --- this
is the part that matters --- \emph{two} qutrits can be correlated in ways that no pair of
classical switches can imitate.
\end{plain}

\begin{center}
\begin{tikzpicture}[font=\small]
  % bit
  \begin{scope}[xshift=0cm]
    \node[font=\bfseries] at (1.1,2.5) {BIT};
    \fill[past!25] (0,0) rectangle (0.9,1.6);
    \fill[past!70] (1.3,0) rectangle (2.2,1.6);
    \node at (0.45,-0.35) {$0$}; \node at (1.75,-0.35) {$1$};
    \node[font=\scriptsize,text=rule,align=center] at (1.1,-1.05)
      {two settings\\one of them at a time};
  \end{scope}
  % trit
  \begin{scope}[xshift=4.4cm]
    \node[font=\bfseries] at (1.5,2.5) {TRIT};
    \fill[future!20] (0,0) rectangle (0.8,1.6);
    \fill[future!45] (1.1,0) rectangle (1.9,1.6);
    \fill[future!75] (2.2,0) rectangle (3.0,1.6);
    \node at (0.4,-0.35) {$0$}; \node at (1.5,-0.35) {$1$}; \node at (2.6,-0.35) {$2$};
    \node[font=\scriptsize,text=rule,align=center] at (1.5,-1.05)
      {three settings\\still one at a time};
  \end{scope}
  % qutrit
  \begin{scope}[xshift=9.6cm]
    \node[font=\bfseries] at (1.5,2.5) {QUTRIT};
    \shade[left color=future!20,right color=future!80] (0,0) rectangle (3.0,1.6);
    \draw[dashed,draw=rule] (1.0,0) -- (1.0,1.6);
    \draw[dashed,draw=rule] (2.0,0) -- (2.0,1.6);
    \node at (0.5,-0.35) {$0$}; \node at (1.5,-0.35) {$1$}; \node at (2.5,-0.35) {$2$};
    \node[font=\scriptsize,text=rule,align=center] at (1.5,-1.05)
      {three settings\\\emph{all at once}, with phases};
  \end{scope}
\end{tikzpicture}
\end{center}

\begin{spec}
A qutrit is a unit vector in $\mathbb{C}^3$. Two qutrits live in
$\mathbb{C}^3\otimes\mathbb{C}^3=\mathbb{C}^9$. The relevant operators are generated by
\[
  X\ket{j}=\ket{j+1 \bmod 3},\qquad
  Z\ket{j}=\omega^{j}\ket{j},\qquad \omega=e^{2\pi i/3},
\]
which satisfy $ZX=\omega XZ$. Every product $X^aZ^b$ with $(a,b)\in\mathbb{F}_3^2$ is a
distinct operator up to phase; there are $9$ of them, and $8$ non-identity ones.
\end{spec}

\subsection{The self-entanglement, and where the $40$ comes from}

\begin{plain}
Now the key move. Take \emph{one} qutrit and entangle it with itself at two different
times --- call them ``past'' and ``future''. Physically this is done by splitting a
single photon's degrees of freedom: its arrival \emph{time} carries one qutrit and its
\emph{colour} carries the other. It is one particle, entangled with its own history.

Mathematically, an operator acting on one qutrit is a $3\times 3$ table of numbers, and
there are $9$ independent slots in such a table. So the space of things you can \emph{do}
to one qutrit is itself $9$-dimensional --- exactly the size of the space of \emph{states}
of two qutrits. This is not a coincidence or an analogy. It is an identity:
\[
  \underbrace{\mathbb{C}^3\otimes\mathbb{C}^3}_{\text{two qutrits}}
  \;=\;
  \underbrace{\mathrm{End}(\mathbb{C}^3)}_{\text{operations on \emph{one} qutrit}}.
\]
So ``the two-qutrit machine'' and ``the one-qutrit machine that can act on itself'' are
the same machine described twice.

Now count the directions. There are $81$ ways to label a pair (four trits), one of which
is ``do nothing''. That leaves $80$. Each direction and its opposite behave identically
for our purposes, and each direction comes in $2$ such flavours, so the distinct
directions number $80/2=40$.
\end{plain}

\headline{$\dfrac{3^4-1}{2}=\dfrac{81-1}{2}=40$. \quad
The forty points of the geometry are the forty independent things one self-entangled
qutrit can be measured along.}

\begin{center}
\begin{tikzpicture}[font=\small,>={Latex[length=2.2mm]}]
  \fill[past!18,rounded corners=3pt] (-3.4,-1.5) rectangle (-0.4,1.5);
  \fill[future!18,rounded corners=3pt] (0.4,-1.5) rectangle (3.4,1.5);
  \node[font=\bfseries,text=past] at (-1.9,1.15) {PAST};
  \node[font=\bfseries,text=future] at (1.9,1.15) {FUTURE};
  \node at (-1.9,0.1) {$\ket{j}_p$};
  \node at (1.9,0.1) {$\ket{j}_f$};
  \node[font=\scriptsize,text=rule] at (-1.9,-0.8) {arrival time};
  \node[font=\scriptsize,text=rule] at (1.9,-0.8) {colour (frequency)};

  \draw[decorate,decoration={snake,amplitude=0.7pt,segment length=4pt},
        line width=1pt,draw=ink] (-0.4,0.1) -- (0.4,0.1);
  \node[font=\scriptsize,align=center] at (0,0.85) {$\ket{\Omega}$};

  \node[align=left,text width=6cm,font=\footnotesize] at (8.4,0.2)
    {$\displaystyle \ket{\Omega}=\mathrm{CX}_{p\to f}\,(F_3\otimes I)\ket{0}_p\ket{0}_f
      =\tfrac{1}{\sqrt3}\sum_{j=0}^{2}\ket{j}_p\ket{j}_f$\\[5pt]
     \textit{One photon. Two degrees of freedom. The entanglement is between the
     particle and itself.}};
\end{tikzpicture}
\end{center}

\begin{gotwrong}
The sentence ``one qutrit, self-entangled as past/future, is enough to generate the full
substrate'' is \emph{not} ours. It is stated verbatim in this project's own
encyclopedia, \cmd{docs/index.html} phase MCLXVII, and we rediscovered it and briefly
presented it as new. What survives as new is narrower: the explicit construction of the
$40$ points as projective left$\times$right multiplication classes on
$\mathrm{End}(\mathbb{C}^3)$, and the identification of the $40$ lines as the maximal
\emph{commuting} subalgebras. The lesson --- search for the \emph{result}, not the topic
--- produced three of the tools in \S\ref{sec:verification}.
\end{gotwrong}

% =====================================================================================
\section{The substrate: $W(3,3)$}\label{sec:geometry}
% =====================================================================================

\begin{plain}
Here is what the $40$ directions look like when you write down which of them are
compatible with which.

Two measurements are \emph{compatible} if you can perform both without one disturbing
the other. Of the $40$ directions, each is compatible with exactly $12$ others. Whenever
two directions are compatible, exactly $2$ further directions are compatible with both.
Whenever two are \emph{in}compatible, exactly $4$ are compatible with both.

Those four numbers --- $40$, $12$, $2$, $4$ --- pin the shape down completely. There is
essentially one object in the world with that compatibility pattern, and it has been
studied since the 1900s under the name $W(3,3)$.
\end{plain}

\begin{spec}
$W(3,3)$ is the symplectic generalised quadrangle of order $(3,3)$: the points are the
$40$ points of $\mathrm{PG}(3,3)$ and the lines are the $40$ totally isotropic lines of a
non-degenerate alternating form. Its collinearity graph is strongly regular with
parameters
\[
  \mathrm{SRG}(40,\,12,\,2,\,4),
\]
i.e.\ $40$ vertices, every vertex of degree $12$, adjacent pairs with $\lambda=2$ common
neighbours, non-adjacent pairs with $\mu=4$. Its automorphism group is
$\mathrm{PGSp}(4,3)=U_4(2){:}2=W(E_6)$ of order $51{,}840$.
\end{spec}

\subsection{Why this shape and not another}

\begin{center}
\begin{tikzpicture}[font=\footnotesize,
  pt/.style={circle,fill=ink,inner sep=1.35pt},
  ln/.style={draw=past,line width=0.7pt,opacity=0.55}]

  % a schematic of the GQ(3,3) incidence: 5 "spread" groups of 8 points on a ring
  \foreach \k in {0,...,7}{
    \pgfmathsetmacro\ang{\k*45}
    \foreach \r/\lab in {2.9/a, 2.15/b, 1.45/c, 0.78/d, 0.28/e}{
      \node[pt] (P\k\lab) at (\ang:\r) {};
    }
  }
  \foreach \k in {0,...,7}{
    \draw[ln] (P\k a) -- (P\k b) -- (P\k c) -- (P\k d) -- (P\k e);
  }
  \foreach \k [evaluate=\k as \j using {int(mod(\k+1,8))}] in {0,...,7}{
    \draw[ln,draw=future] (P\k a) -- (P\j b);
    \draw[ln,draw=future,opacity=0.3] (P\k c) -- (P\j d);
  }
  \node[align=center,text width=5.4cm] at (7.4,1.2)
   {\textbf{$40$ points, $40$ lines.}\\[3pt]
    Every point lies on exactly $4$ lines.\\
    Every line carries exactly $4$ points.\\
    Given a line $\ell$ and a point $P\notin\ell$,\\
    there is \emph{exactly one} point of $\ell$\\
    collinear with $P$.};
  \node[align=center,text width=5.4cm,font=\scriptsize,text=rule] at (7.4,-1.6)
   {\textit{(Schematic. The full incidence graph has $160$ flags and does not embed in
    the plane without crossings; this figure shows the ring structure and two of the
    line families only.)}};
\end{tikzpicture}
\end{center}

\begin{plain}
That last property --- ``exactly one'' --- is the entire content of the word
\emph{quadrangle}. It is a strong constraint: it means the geometry contains no
triangles, so there is never any ambiguity about how to get from one place to another.
For a computer, that is precisely what you want from an address space and from a routing
fabric: there is one shortest path, and no loops to break ties in.
\end{plain}

\subsection{The two groups of order $51{,}840$, which are not the same group}

\begin{spec}
Two distinct groups of order $51{,}840$ appear in this project, and confusing them is the
most common error in this area:
\[
  \mathrm{Sp}(4,3)=2.U_4(2)\quad(\text{centre of order }2),\qquad
  \mathrm{PGSp}(4,3)=U_4(2){:}2=W(E_6)\quad(\text{trivial centre}).
\]
They are \emph{not} isomorphic. The first is the group acting on Pauli frames --- the one
the hardware implements. The second is the automorphism group of the geometry --- the one
the drawing has. One is a central double cover; the other is an outer extension.
\end{spec}

\begin{plain}
This distinction is not pedantry, and it has physical content. The extra element in the
first group acts as a sign, $z\mapsto -1$. Whether that sign is ``spent'' or ``free''
decides whether the resulting structure has a handedness --- whether it distinguishes
left from right. This is why the machine can \emph{host} a chirality, though as we
established in earlier work it cannot \emph{explain} one from the inside.
\end{plain}

% =====================================================================================
